% ============================================================
% HSBI Beamer — ISLP Lecture Series
% Final Mock Exam — Set C — Solutions (Lectures 1-12, comprehensive)
% Compile: pdflatex -jobname=Final_Mock_Exam_C_Solutions_Slides solutions_slides_c.tex
% ============================================================
\documentclass[aspectratio=169,11pt]{beamer}
\usepackage[utf8]{inputenc}\usepackage[T1]{fontenc}\usepackage[english]{babel}
\usepackage{graphicx}\usepackage{booktabs}\usepackage{amsmath,amssymb,amsfonts}
\usepackage{mathtools}\usepackage{hyperref}\usepackage{xcolor}
\usepackage{verbatim}\usepackage{alltt}
\usepackage{tikz}\usetikzlibrary{calc,arrows.meta,positioning,shapes}
\usepackage{tcolorbox}\tcbuselibrary{skins}

\usetheme{Madrid}\usecolortheme{default}
\setbeamertemplate{navigation symbols}{}
\setbeamertemplate{footline}[frame number]
\setbeamertemplate{blocks}[rounded][shadow=false]
\setbeamertemplate{headline}{%
  \leavevmode\hbox{%
    \begin{beamercolorbox}[wd=\paperwidth,ht=2.2ex,dp=0.9ex]{section in head/foot}%
      {\fontsize{4.6}{5.2}\selectfont%
       \insertsectionnavigationhorizontal{\paperwidth}{}{\hskip0pt plus1filll}}%
    \end{beamercolorbox}}%
}
\setbeamerfont{section in head/foot}{size=\tiny}
\setbeamerfont{palette primary}{size=\tiny}
\setbeamerfont{palette tertiary}{size=\tiny}
\setbeamerfont{section in toc}{size=\footnotesize}

\usepackage{listings}
\definecolor{pyKw}{RGB}{0,0,180}\definecolor{pyStr}{RGB}{20,120,40}
\definecolor{pyCom}{RGB}{120,120,120}\definecolor{accent}{RGB}{38,70,140}
\lstdefinestyle{Pystyle}{language=Python,basicstyle=\ttfamily\footnotesize,
  keywordstyle=\color{pyKw}\bfseries,commentstyle=\color{pyCom}\itshape,
  stringstyle=\color{pyStr},showstringspaces=false,upquote=true,
  columns=fullflexible,breaklines=true,literate={~}{{$\sim$}}1}
\lstset{style=Pystyle}

\newtcolorbox{takeaway}[1][Takeaway]{enhanced,colback=green!4,colframe=green!50!black,
  boxrule=0pt,leftrule=2.5pt,arc=2pt,fonttitle=\bfseries,title=#1,
  left=2mm,right=2mm,top=1mm,bottom=1mm}
\newtcolorbox{numexample}[1][Worked example]{enhanced,colback=orange!4,colframe=orange!70!black,
  boxrule=0pt,leftrule=2.5pt,arc=2pt,fonttitle=\bfseries,title=#1,
  left=2mm,right=2mm,top=1mm,bottom=1mm}
\newtcolorbox{readme}[1][How to read this]{enhanced,colback=blue!4,colframe=accent,
  boxrule=0pt,leftrule=2.5pt,arc=2pt,fonttitle=\bfseries,title=#1,
  left=2mm,right=2mm,top=1mm,bottom=1mm}

\title{Final Mock Exam --- Set C --- Solutions}
\subtitle{Lectures 1--12 (comprehensive)}
\author{Prof.\ Dr.\ Christoph Weisser}\institute{Quantitative Research Methods \textperiodcentered\ HSBI}
\date{Summer Semester 2026}\graphicspath{{images/}}

\newtcolorbox{exercise}[1][Exercise]{enhanced, fontupper=\footnotesize, colback=purple!4, colframe=purple!55!black, boxrule=0pt, leftrule=2.5pt, arc=2pt, fonttitle=\bfseries, title=#1, left=2mm, right=2mm, top=1mm, bottom=1mm}
\newtcolorbox{solutionbox}[1][Solution]{enhanced, fontupper=\footnotesize, colback=teal!5, colframe=teal!55!black, boxrule=0pt, leftrule=2.5pt, arc=2pt, fonttitle=\bfseries, title=#1, left=2mm, right=2mm, top=1mm, bottom=1mm}

% Reclaim ~2pt of body height (custom nav headline makes full slides tip over by ~1.83pt)
\addtobeamertemplate{frametitle}{}{\vspace{-2pt}}

\begin{document}
\begin{frame}[plain]\titlepage\end{frame}

% ------------------------------------------------------------
\begin{frame}{Overview of the final mock exam --- Set C}
  \scriptsize
  \begin{center}
  \begin{tabular}{@{}cllc@{}}
    \toprule
    \textbf{Problem} & \textbf{Topic} & \textbf{Lecture slides} & \textbf{Points}\\
    \midrule
    P1 & Moving beyond linearity --- splines & Ch.\ 7 --- Moving Beyond Linearity (Lecture 9) & 16\\
    P2 & Generalised additive models & Ch.\ 7 --- Moving Beyond Linearity (Lecture 9) & 8\\
    P3 & Trees by hand (regression \& classification) & Ch.\ 8 --- Tree-Based Methods (Lecture 10) & 20\\
    P4 & Ensembles --- bagging, RF, boosting & Ch.\ 8 --- Tree-Based Methods (Lecture 10) & 16\\
    P5 & Neural networks by hand & Ch.\ 10 --- Deep Learning (Lecture 11) & 20\\
    P6 & Multiple testing & Ch.\ 13 --- Multiple Testing (Lecture 12) & 20\\
    P7 & Cumulative essentials & Ch.\ 2--6 (Lectures 1--8) & 20\\
    \midrule
    $\Sigma$ & & & \textbf{120}\\
    \bottomrule
  \end{tabular}
  \end{center}
  \vspace{2pt}
  \begin{takeaway}[Exam conditions]
    120 minutes for 120 points ($\approx$ 1 point per minute). Allowed aids: non-programmable
    calculator and one handwritten A4 sheet (both sides). All required constants ($e$-powers
    and logs) are given in the problems. Justify every answer.
  \end{takeaway}
\end{frame}

% ============================================================
% P1
% ============================================================
\begin{frame}{Problem 1 --- task}
  \begin{exercise}[Problem 1 (16 P) --- Moving beyond linearity]
    {\scriptsize\itshape Lecture slides: Chapter 7 --- Moving Beyond Linearity (Lecture 9).}\par\smallskip
    \begin{itemize}
      \item[(a)] \textbf{[4 P]} Degrees of freedom of a cubic spline with $K$ interior knots ---
        derive by counting parameters and constraints; evaluate for $K=3$.
      \item[(b)] \textbf{[4 P]} Degrees of freedom of a \emph{natural} cubic spline with the same
        $K=3$ knots; which extra constraint acts where, and why is it useful?
      \item[(c)] \textbf{[4 P]} Smoothing spline
        $\sum_i (y_i-g(x_i))^2+\lambda\int g''(t)^2dt$: role of the two terms;
        behaviour of $\hat g$ as $\lambda\to 0$ and $\lambda\to\infty$.
      \item[(d)] \textbf{[4 P]} Truncated-power basis of a cubic spline with one knot $\xi$;
        define $(x-\xi)^3_+$; why does adding it keep the fit smooth at $\xi$?
    \end{itemize}
  \end{exercise}
\end{frame}

\begin{frame}{Problem 1 --- solution (a) and (b)}
  \begin{solutionbox}[Solution P1 --- (a) and (b)]
    \textbf{(a)} $K$ knots $\Rightarrow$ $K+1$ cubic pieces $\times\,4$ coefficients $=4K+4$
    parameters; continuity of $g$ and $g'$ and $g''$ at each knot $\Rightarrow$ $3K$ constraints:
    \[
      \mathrm{df}=4(K+1)-3K=K+4
      \qquad\Rightarrow\qquad K=3:\ \mathrm{df}=\mathbf{7}.
    \]
    (Equivalently: basis $1$, $x$, $x^2$, $x^3$ plus one truncated power per knot.)
    \medskip

    \textbf{(b)} Natural spline: \emph{linear beyond the boundary knots}
    (regions $X<\xi_1$ and $X>\xi_K$); killing the quadratic and cubic terms there costs
    $2$ parameters per boundary region:
    \[
      \mathrm{df}=(K+4)-4=K
      \qquad\Rightarrow\qquad K=3:\ \mathrm{df}=\mathbf{3}.
    \]
    \emph{Why useful:} splines are most erratic near the boundaries where data are sparse;
    the linearity constraint stabilises the fit there (lower variance for a little extra bias).
  \end{solutionbox}
\end{frame}

\begin{frame}{Problem 1 --- solution (c) and (d)}
  \begin{solutionbox}[Solution P1 --- (c) and (d)]
    \textbf{(c)} RSS term $=$ fidelity to the data; $\lambda\int g''(t)^2dt$ penalises curvature
    (wiggliness); $\lambda$ trades the two off.
    \begin{itemize}
      \item $\lambda\to 0$: no penalty $\Rightarrow$ $\hat g$ \emph{interpolates} the training
        points exactly --- extremely wiggly, zero training RSS, severe overfitting
        (effective df $\to n$).
      \item $\lambda\to\infty$: curvature infinitely expensive $\Rightarrow$ $g''\equiv 0$
        $\Rightarrow$ $\hat g$ is the \emph{least-squares straight line} (effective df $\to 2$).
    \end{itemize}
    $\lambda$ (chosen by LOOCV) steers the effective df smoothly between $n$ and $2$.
    \medskip

    \textbf{(d)} $f(x)=\beta_0+\beta_1 x+\beta_2 x^2+\beta_3 x^3+\beta_4 (x-\xi)^3_+$ with
    \[
      (x-\xi)^3_+=\begin{cases}(x-\xi)^3 & x>\xi\\ 0 & x\le\xi\end{cases}
    \]
    $(x-\xi)^3_+$ has value and first and second derivative $0$ at $\xi$ $\Rightarrow$ only the
    third derivative jumps: the fit stays continuous with continuous $g'$ and $g''$ --- a cubic spline.
  \end{solutionbox}
\end{frame}

% ============================================================
% P2
% ============================================================
\begin{frame}{Problem 2 --- task}
  \begin{exercise}[Problem 2 (8 P) --- Generalised additive models]
    {\scriptsize\itshape Lecture slides: Chapter 7 --- Moving Beyond Linearity (Lecture 9).}\par\smallskip
    \begin{itemize}
      \item[(a)] \textbf{[3 P]} General form of a GAM for quantitative $Y$; how does it generalise
        multiple linear regression?
      \item[(b)] \textbf{[3 P]} One advantage concerning the interpretation of the $f_j$;
        the main structural limitation --- and its repair.
      \item[(c)] \textbf{[2 P]} Backfitting in one or two sentences.
    \end{itemize}
  \end{exercise}
\end{frame}

\begin{frame}{Problem 2 --- solution}
  \begin{solutionbox}[Solution P2 --- GAMs]
    \textbf{(a)}\ \ $Y=\beta_0+f_1(X_1)+f_2(X_2)+\dots+f_p(X_p)+\varepsilon$.\\
    Linear regression is the special case $f_j(X_j)=\beta_j X_j$; a GAM replaces each linear
    term by a smooth non-linear $f_j$ (spline, local regression) but keeps the
    \emph{additive} structure.
    \medskip

    \textbf{(b)} \emph{Advantage:} additivity $\Rightarrow$ the effect of each $X_j$ is its own
    function $f_j$, which can be plotted and interpreted individually while the other predictors
    are held fixed --- one panel per predictor, as readable as a coefficient table.\\
    \emph{Limitation:} no \emph{interactions} --- the effect of $X_j$ cannot depend on $X_k$.\\
    \emph{Repair:} add interaction components manually, e.g.\ $X_j\times X_k$ or a
    two-dimensional smooth $f_{jk}(X_j{,}X_k)$.
    \medskip

    \textbf{(c)} \emph{Backfitting:} cycle through $j=1{,}\dots{,}p$ and re-fit each $f_j$ to the
    partial residuals $r_i=y_i-\beta_0-\sum_{k\ne j}f_k(x_{ik})$ with all other functions held
    fixed; repeat until convergence.
  \end{solutionbox}
\end{frame}

% ============================================================
% P3
% ============================================================
\begin{frame}{Problem 3 --- task}
  \begin{exercise}[Problem 3 (20 P) --- Trees by hand]
    {\scriptsize\itshape Lecture slides: Chapter 8 --- Tree-Based Methods (Lecture 10).}\par\smallskip
    Delivery vans $x$ and parcels delivered $y$ (1{,}000) for $n=6$ depots;
    one split at $s$: $R_1=\{x<s\}$ and $R_2=\{x\ge s\}$ with region-mean predictions.
    \begin{center}\scriptsize
    \begin{tabular}{lcccccc}
      \toprule
      $i$ & 1 & 2 & 3 & 4 & 5 & 6\\
      \midrule
      $x_i$ (vans) & 1 & 2 & 3 & 4 & 5 & 6\\
      $y_i$ (parcels) & 3 & 5 & 4 & 9 & 11 & 10\\
      \bottomrule
    \end{tabular}
    \end{center}
    (a) \textbf{[8 P]} Evaluate the cutpoints $s=2.5$ and $s=3.5$ by total RSS; which is chosen?
    (b) \textbf{[2 P]} Prediction for a new depot with $x=5$ vans.
    (c) \textbf{[6 P]} Node with 5 ``Electronics'' / 3 ``Clothing'' / 1 ``Books'': proportions;
        Gini; entropy ($\ln\tfrac59=-0.5878$; $\ln\tfrac13=-1.0986$; $\ln\tfrac19=-2.1972$);
        predicted class.
    (d) \textbf{[4 P]} Cost-complexity pruning: role of $\alpha$ in
        $\sum_m\sum_{i\in R_m}(y_i-\hat y_{R_m})^2+\alpha|T|$.
  \end{exercise}
\end{frame}

\begin{frame}{Problem 3 --- solution (a) and (b)}
  \begin{solutionbox}[Solution P3 --- (a) and (b)]
    \textbf{(a)} Region means and RSS per cutpoint ($\bar y=42/6=7$; root RSS $=58$):
    \begin{center}\scriptsize
    \begin{tabular}{lccccc}
      \toprule
      $s$ & $R_1$ ($y$-values) & $\hat y_{R_1}$ & $R_2$ ($y$-values) & $\hat y_{R_2}$ & total RSS\\
      \midrule
      2.5 & $\{3{,}5\}$ & 4 & $\{4{,}9{,}11{,}10\}$ & 8.5 & $2+29=\mathbf{31}$\\
      3.5 & $\{3{,}5{,}4\}$ & 4 & $\{9{,}11{,}10\}$ & 10 & $2+2=\mathbf{4}$\\
      \bottomrule
    \end{tabular}
    \end{center}
    Details: $s=2.5$: $(3-4)^2+(5-4)^2=2$; $(4-8.5)^2+(9-8.5)^2+(11-8.5)^2+(10-8.5)^2=29$.\\
    $s=3.5$: $(3-4)^2+(5-4)^2+(4-4)^2=2$; $(9-10)^2+(11-10)^2+(10-10)^2=2$.\\
    Greedy recursive binary splitting picks the smaller RSS: $\mathbf{s=3.5}$ ($4\ll 31$; in fact
    the best of all five cutpoints) --- it isolates the three high-volume depots.
    \medskip

    \textbf{(b)} $x=5\ge 3.5$ $\Rightarrow$ region $R_2$ $\Rightarrow$
    $\hat y=\hat y_{R_2}=\mathbf{10}$ (i.e.\ 10{,}000 parcels).
  \end{solutionbox}
\end{frame}

\begin{frame}{Problem 3 --- solution (c) and (d)}
  \begin{solutionbox}[Solution P3 --- (c) and (d)]
    \textbf{(c)} $\hat p_{\text{E}}=5/9$; $\hat p_{\text{C}}=3/9=1/3$; $\hat p_{\text{B}}=1/9$.
    \[
      G=\tfrac59\cdot\tfrac49+\tfrac39\cdot\tfrac69+\tfrac19\cdot\tfrac89
       =\tfrac{20+18+8}{81}=\tfrac{46}{81}=\mathbf{0.568}
    \]
    \[
      D=-\bigl(\tfrac59\ln\tfrac59+\tfrac13\ln\tfrac13+\tfrac19\ln\tfrac19\bigr)
       =0.3266+0.3662+0.2441=\mathbf{0.937}
    \]
    Predicted class: majority $=$ \textbf{Electronics} (estimated probability $5/9\approx 0.556$).
    Both measures are $0$ for a pure node --- this node is quite impure (no class dominates).
    \medskip

    \textbf{(d)} $\alpha$ $=$ price per terminal node $|T|$:
    \begin{itemize}
      \item $\alpha=0$: criterion $=$ training RSS $\Rightarrow$ full tree $T_0$ is optimal.
      \item $\alpha\uparrow$: leaves must pay for themselves; branches with small RSS gain are
        pruned $\Rightarrow$ nested sequence of ever smaller subtrees (weakest-link pruning).
      \item Choice: \emph{cross-validation} over $\alpha$; refit the chosen subtree on all data
        --- balances bias against variance.
    \end{itemize}
  \end{solutionbox}
\end{frame}

% ============================================================
% P4
% ============================================================
\begin{frame}{Problem 4 --- task}
  \begin{exercise}[Problem 4 (16 P) --- Ensembles]
    {\scriptsize\itshape Lecture slides: Chapter 8 --- Tree-Based Methods (Lecture 10).}\par\smallskip
    \begin{itemize}
      \item[(a)] \textbf{[6 P]} $B$ identically distributed trees with variance $\sigma^2$ and
        pairwise correlation $\rho$:
        $\operatorname{Var}\bigl(\frac1B\sum_b\hat f_b(x)\bigr)=\rho\sigma^2+\frac{1-\rho}{B}\sigma^2$.
        (i) Why does this show variance reduction and what limits it?
        (ii) Evaluate for $B=50$ and $\rho=0.4$ and $\sigma^2=16$; compare with one tree.
        (iii) Limit as $B\to\infty$?
      \item[(b)] \textbf{[4 P]} Random forests: modification at each split; recommended $m$ for
        \emph{regression} with $p=30$; why does it improve on bagging (use the formula)?
      \item[(c)] \textbf{[4 P]} Boosting: the three tuning parameters; the slow-learning idea.
      \item[(d)] \textbf{[2 P]} What does the OOB error estimate and why is it free?
    \end{itemize}
  \end{exercise}
\end{frame}

\begin{frame}{Problem 4 --- solution (a) and (b)}
  \begin{solutionbox}[Solution P4 --- (a) and (b)]
    \textbf{(a)} (i) Averaging drives the second term $\frac{1-\rho}{B}\sigma^2\to 0$
    (the independent part of the trees' variability); the first term $\rho\sigma^2$ is
    independent of $B$ --- correlation between bootstrap trees limits the reduction.\\
    (ii) $\operatorname{Var}=0.4\cdot 16+\frac{0.6}{50}\cdot 16=6.4+0.192=\mathbf{6.592}$
    vs.\ $\sigma^2=16$ for a single tree ($\approx 58.8\%$ less).\\
    (iii) $B\to\infty$: variance $\to\rho\sigma^2=\mathbf{6.4}$ (floor; more trees never
    overfit, they only stabilise).
    \medskip

    \textbf{(b)} At \emph{each split} only a random subset of $m$ of the $p$ predictors may be
    used; regression default $m\approx p/3$ $\Rightarrow$ $p=30$: $m=30/3=\mathbf{10}$.\\
    With bagging a dominant predictor tops nearly every tree $\Rightarrow$ trees alike
    $\Rightarrow$ $\rho$ high. Restricting candidates \emph{decorrelates} the trees
    ($\rho\downarrow$), which lowers the floor $\rho\sigma^2$ --- the part of the variance
    that averaging alone cannot remove.
  \end{solutionbox}
\end{frame}

\begin{frame}{Problem 4 --- solution (c) and (d)}
  \begin{solutionbox}[Solution P4 --- (c) and (d)]
    \textbf{(c)} Boosting for regression trees:
    \begin{itemize}
      \item \textbf{$B$ --- number of trees:} added sequentially; too many $\Rightarrow$ slow
        overfitting; choose by CV.
      \item \textbf{$\lambda$ --- shrinkage} (e.g.\ $0.01$ or $0.001$): scales each new tree's
        contribution; smaller $\lambda$ $\Rightarrow$ slower learning $\Rightarrow$ larger $B$ needed.
      \item \textbf{$d$ --- splits per tree} (interaction depth): complexity and interaction
        order of each tree; often stumps ($d=1$) suffice (additive model).
    \end{itemize}
    \emph{Slow learning:} each tree is fitted to the current \emph{residuals} and only the
    fraction $\lambda$ is added --- the model improves gradually exactly where it still errs,
    instead of fitting the data hard at once; this tends to generalise better.
    \medskip

    \textbf{(d)} Each bootstrap sample misses $\approx 1/3$ of the observations (out-of-bag);
    predicting every observation only with trees that never saw it yields the \textbf{OOB error}
    --- a valid \emph{test-error} estimate, free as a by-product of the bootstrap
    (no validation set or CV loop needed).
  \end{solutionbox}
\end{frame}

% ============================================================
% P5
% ============================================================
\begin{frame}{Problem 5 --- task}
  \begin{exercise}[Problem 5 (20 P) --- Neural networks by hand]
    {\scriptsize\itshape Lecture slides: Chapter 10 --- Deep Learning (Lecture 11).}\par\smallskip
    Input $x=(2{,}3)$; hidden ReLU units $g(z)=\max\{0{,}z\}$; linear output:
    \[
      A_1=g(-3+2x_1+1x_2)\qquad A_2=g(2-2x_1-1x_2)\qquad f(x)=-2+3A_1+2A_2
    \]
    \begin{itemize}
      \item[(a)] \textbf{[8 P]} Forward pass step by step: pre-activations; $A_1$ and $A_2$;
        output $f(x)$. What did the ReLU do to unit 2 and why are non-linear activations essential?
      \item[(b)] \textbf{[6 P]} Logits $z=(2.0{,}\ 0.0{,}\ -1.0)$; true class 2: all three softmax
        probabilities and the cross-entropy loss
        ($e^{2}=7.389$; $e^{0}=1.000$; $e^{-1}=0.368$; $\ln 8.757=2.170$).
      \item[(c)] \textbf{[4 P]} Parameter count of a $p=6\to k=5\to K=3$ network with biases
        (show both layers).
      \item[(d)] \textbf{[2 P]} Conv layer: $64\times 64$ input; $F=3$; $P=1$; $S=2$:
        output size via $\lfloor(W-F+2P)/S\rfloor+1$.
    \end{itemize}
  \end{exercise}
\end{frame}

\begin{frame}{Problem 5 --- solution (a) forward pass}
  \begin{solutionbox}[Solution P5 --- (a) forward pass]
    Pre-activations:
    \[
      z_1=-3+2\cdot 2+1\cdot 3=4
      \qquad\qquad
      z_2=2-2\cdot 2-1\cdot 3=-5
    \]
    ReLU activations:
    \[
      A_1=\max\{0{,}4\}=4\qquad\qquad A_2=\max\{0{,}-5\}=0
    \]
    Output:
    \[
      f(x)=-2+3\cdot 4+2\cdot 0=\mathbf{10}
    \]
    \medskip
    The ReLU \emph{switched unit 2 off}: its negative pre-activation is clipped to $0$, so it
    contributes nothing for this input (though it can for others).\\
    \emph{Why essential:} with a linear $g$ the whole network collapses into one linear function
    of $x$ regardless of depth --- the non-linearity is what lets hidden layers build flexible
    functions (piecewise-linear surfaces with ReLU).
  \end{solutionbox}
\end{frame}

\begin{frame}{Problem 5 --- solution (b) softmax and cross-entropy}
  \begin{solutionbox}[Solution P5 --- (b) softmax and cross-entropy]
    Exponentials and their sum:
    \[
      e^{2.0}=7.389\qquad e^{0.0}=1.000\qquad e^{-1.0}=0.368
      \qquad\Rightarrow\qquad 7.389+1.000+0.368=8.757
    \]
    Softmax probabilities:
    \[
      \hat p_1=\frac{7.389}{8.757}=\mathbf{0.844}\qquad
      \hat p_2=\frac{1.000}{8.757}=\mathbf{0.114}\qquad
      \hat p_3=\frac{0.368}{8.757}=\mathbf{0.042}
    \]
    (sum $=1.000$ \checkmark). Cross-entropy for the \emph{true} class 2:
    \[
      -\ln\hat p_2=-\ln\frac{1.000}{8.757}=-(0-\ln 8.757)=\ln 8.757=\mathbf{2.170}
    \]
    The model's most probable class is class 1 ($\hat p_1=0.844$), but the loss is evaluated at
    the true class 2 (only $0.114$) --- hence the sizeable loss; a perfect $\hat p_2=1$ gives $0$.
  \end{solutionbox}
\end{frame}

\begin{frame}{Problem 5 --- solution (c) and (d)}
  \begin{solutionbox}[Solution P5 --- (c) parameters and (d) convolution]
    \textbf{(c)} Fully connected $6\to 5\to 3$ with biases:
    \begin{center}
    \begin{tabular}{@{}lccc@{}}
      \toprule
      Layer & Weights & Biases & Parameters\\
      \midrule
      input $\to$ hidden & $6\cdot 5=30$ & $5$ & $35$\\
      hidden $\to$ output & $5\cdot 3=15$ & $3$ & $18$\\
      \midrule
      total & & & $\mathbf{53}$\\
      \bottomrule
    \end{tabular}
    \end{center}
    \medskip

    \textbf{(d)} Output width
    $=\Bigl\lfloor\dfrac{W-F+2P}{S}\Bigr\rfloor+1
      =\Bigl\lfloor\dfrac{64-3+2\cdot 1}{2}\Bigr\rfloor+1
      =\lfloor 31.5\rfloor+1=32$
    $\Rightarrow$ output feature map $\mathbf{32\times 32}$: stride $S=2$ roughly halves the
    $64\times 64$ input (the floor discards the fractional part before adding $1$).
  \end{solutionbox}
\end{frame}

% ============================================================
% P6
% ============================================================
\begin{frame}{Problem 6 --- task}
  \begin{exercise}[Problem 6 (20 P) --- Multiple testing]
    {\scriptsize\itshape Lecture slides: Chapter 13 --- Multiple Testing (Lecture 12).}\par\smallskip
    \begin{itemize}
      \item[(a)] \textbf{[4 P]} $m=40$ independent tests at $\alpha=0.05$; all nulls true:
        define and compute the FWER ($(0.95)^{40}=0.1285$); interpret.
      \item[(b)] \textbf{[7 P]} Six ordered $p$-values
        $0.001$; $0.008$; $0.011$; $0.030$; $0.041$; $0.260$ --- FWER control at $\alpha=0.05$
        with (i) Bonferroni and (ii) Holm; thresholds; rejected hypotheses; why is Holm never
        weaker than Bonferroni?
      \item[(c)] \textbf{[6 P]} Benjamini--Hochberg at $q=0.05$ on the same $p$-values:
        thresholds $jq/m$; rejections; what exactly does FDR control guarantee?
      \item[(d)] \textbf{[3 P]} Confirmatory phase-III trial with $m=5$ pre-specified endpoints
        (one false positive $\Rightarrow$ an ineffective drug approved): FWER or FDR? Justify.
    \end{itemize}
  \end{exercise}
\end{frame}

\begin{frame}{Problem 6 --- solution (a) FWER}
  \begin{solutionbox}[Solution P6 --- (a) FWER for 40 independent tests]
    $\text{FWER}=\Pr(V\ge 1)$ --- the probability of \emph{at least one} false rejection among
    the $m$ tests ($V$ $=$ number of falsely rejected true nulls).
    \medskip

    With all 40 nulls true and independent tests at level $0.05$:
    \[
      \text{FWER}=1-\Pr(\text{no false rejection})
      =1-(1-0.05)^{40}=1-(0.95)^{40}=1-0.1285=\mathbf{0.8715}
    \]
    \medskip

    \emph{Interpretation:} although each test individually errs with probability $5\%$ only,
    the chance of at least one spurious ``discovery'' among 40 unadjusted tests is about
    $87\%$ --- multiple testing without adjustment almost guarantees false findings.
  \end{solutionbox}
\end{frame}

\begin{frame}{Problem 6 --- solution (b) Bonferroni and Holm}
  \begin{solutionbox}[Solution P6 --- (b) Bonferroni and Holm at level 0.05]
    \textbf{Bonferroni:} reject iff $p_{(j)}\le\alpha/m=0.05/6=0.00833$:
    only $0.001$ and $0.008$ pass ($0.011>0.00833$) $\Rightarrow$
    reject $H_{(1)}$ and $H_{(2)}$: \textbf{2 rejections}.
    \medskip

    \textbf{Holm (step-down):} compare $p_{(j)}$ with $\alpha/(m-j+1)$; stop at first failure:
    \begin{center}\scriptsize
    \begin{tabular}{lcccccc}
      \toprule
      $j$ & 1 & 2 & 3 & 4 & 5 & 6\\
      \midrule
      $p_{(j)}$ & 0.001 & 0.008 & 0.011 & 0.030 & 0.041 & 0.260\\
      $\alpha/(m-j+1)$ & 0.00833 & 0.0100 & 0.0125 & 0.0167 & 0.0250 & 0.0500\\
      decision & \checkmark & \checkmark & \checkmark & \textbf{stop} & --- & ---\\
      \bottomrule
    \end{tabular}
    \end{center}
    $0.011\le 0.0125$ but $0.030>0.0167$ $\Rightarrow$ reject
    $H_{(1)}$ and $H_{(2)}$ and $H_{(3)}$: \textbf{3 rejections}.
    \medskip

    \emph{Holm $\supseteq$ Bonferroni:} Holm's hurdles $\alpha/(m-j+1)\ge\alpha/m$ for all $j$
    (equality only at $j=1$) $\Rightarrow$ uniformly more powerful --- and it still controls the
    FWER without independence assumptions. Here Holm rejects one more ($H_{(3)}$).
  \end{solutionbox}
\end{frame}

\begin{frame}{Problem 6 --- solution (c) and (d)}
  \begin{solutionbox}[Solution P6 --- (c) Benjamini--Hochberg and (d) FDR vs FWER]
    \textbf{(c)} Thresholds $jq/m=j\cdot 0.05/6$:
    \begin{center}\scriptsize
    \begin{tabular}{lcccccc}
      \toprule
      $j$ & 1 & 2 & 3 & 4 & 5 & 6\\
      \midrule
      $p_{(j)}$ & 0.001 & 0.008 & 0.011 & 0.030 & 0.041 & 0.260\\
      $jq/m$ & 0.00833 & 0.0167 & 0.0250 & 0.0333 & 0.0417 & 0.0500\\
      $p_{(j)}\le jq/m$? & yes & yes & yes & yes & \textbf{yes} & no\\
      \bottomrule
    \end{tabular}
    \end{center}
    Largest such $j$: $L=5$ ($0.041\le 0.0417$; note $0.260>0.0500$) $\Rightarrow$ reject
    \emph{all} of $H_{(1)}$ to $H_{(5)}$: \textbf{5 rejections} (most liberal: $5$ vs.\ Holm $3$,
    Bonferroni $2$).
    \emph{Guarantee:} $\text{FDR}=E[V/\max(R{,}1)]\le q$ --- \emph{on average} at most $5\%$
    of rejected hypotheses are false discoveries (no bound on $\Pr(V\ge 1)$; that is FWER).
    \medskip

    \textbf{(d)} Control the \textbf{FWER}. A confirmatory trial with only $m=5$ pre-specified
    endpoints, where a single false positive is very costly (an ineffective drug approved),
    demands a bound on the probability of \emph{any} false rejection --- exactly FWER control
    (e.g.\ Holm at $0.05$). FDR only limits the \emph{expected proportion} of false endpoints
    and would tolerate an occasional false positive: fine for exploratory screens, not here.
  \end{solutionbox}
\end{frame}

% ============================================================
% P7
% ============================================================
\begin{frame}{Problem 7 --- task}
  \begin{exercise}[Problem 7 (20 P) --- Cumulative essentials]
    {\scriptsize\itshape Lecture slides: Chapters 2--6 (Lectures 1--8).}\par\smallskip
    \scriptsize
    \begin{itemize}
      \item[(a)] \textbf{[4 P]} Bootstrap: (i) show $\Pr(\text{obs.\ not in sample})=(1-1/n)^n$
        and its limit ($e^{-1}=0.368$); (ii) fraction included as $n\to\infty$;
        (iii) exact inclusion probability for $n=5$ ($0.8^5=0.328$).
      \item[(b)] \textbf{[4 P]} Bias--variance at $x_0$: method (1) $\operatorname{Var}=3.5$,
        $\operatorname{Bias}=-1.0$, $\operatorname{Var}(\varepsilon)=2.0$; method (2)
        $\operatorname{Var}=0.5$, $\operatorname{Bias}=-2.5$, $\operatorname{Var}(\varepsilon)=2.0$.
        Expected test MSE of each; which to deploy; lowest achievable MSE?
      \item[(c)] \textbf{[3 P]} Elastic-net penalty: what does it inherit from lasso and from
        ridge? When does it clearly beat the lasso?
      \item[(d)] \textbf{[3 P]} Collinearity (corr.\ $0.95$): effect on OLS estimates;
        $\mathrm{VIF}_j=1/(1-R_j^2)$ for $R_j^2=0.9$ ($\sqrt{10}=3.162$).
      \item[(e)] \textbf{[6 P]} Match to \{linear regression; lasso; logistic regression;
        random forest\}, each once: (i) $\Delta$wage per school-year with CI ($n=3{,}000$);
        (ii) biomarker from $8{,}000$ metabolites ($n=80$) with a short list;
        (iii) churn yes/no with odds for management;
        (iv) electricity demand from hundreds of interacting variables, accuracy only.
    \end{itemize}
  \end{exercise}
\end{frame}

\begin{frame}{Problem 7 --- solution (a) to (c)}
  \begin{solutionbox}[Solution P7 --- (a) to (c)]
    \textbf{(a)} Each draw misses a fixed observation with prob.\ $1-\tfrac1n$; all $n$ draws
    (independent) miss it with prob.\ $(1-\tfrac1n)^n\to e^{-1}=\mathbf{0.368}$. So it is
    \emph{included} with prob.\ $1-0.368=\mathbf{0.632}$ ($\approx 63.2\%$ of observations
    appear; the rest are OOB). For $n=5$: $1-0.8^5=1-0.328=\mathbf{0.672}$.
    \medskip

    \textbf{(b)} $\text{MSE}=\operatorname{Var}+\operatorname{Bias}^2+\operatorname{Var}(\varepsilon)$:
    \[
      \text{MSE}_1=3.5+(-1.0)^2+2.0=\mathbf{6.5}\qquad
      \text{MSE}_2=0.5+(-2.5)^2+2.0=\mathbf{8.75}
    \]
    Deploy the \textbf{flexible spline} (method 1, $6.5<8.75$): its higher variance is more than
    offset by its much smaller squared bias. Lowest achievable MSE $=$ irreducible error
    $\operatorname{Var}(\varepsilon)=\mathbf{2.0}$.
    \medskip

    \textbf{(c)} Elastic net: $\lambda\bigl[\tfrac{1-\alpha}{2}\sum_j\beta_j^2+\alpha\sum_j|\beta_j|\bigr]$.
    From \emph{lasso} ($\ell_1$): exact zeros / variable selection; from \emph{ridge} ($\ell_2$):
    graceful handling of correlated predictors and selecting $>n$ variables when $p>n$. It beats
    the lasso with \emph{groups of correlated predictors} (lasso picks one and drops the rest;
    elastic net keeps the group --- ``grouping effect'').
  \end{solutionbox}
\end{frame}

\begin{frame}{Problem 7 --- solution (d) and (e)}
  \begin{solutionbox}[Solution P7 --- (d) and (e)]
    \textbf{(d)} Collinearity does not bias OLS but makes it \emph{imprecise/unstable}: inflated
    variances and standard errors, wide CIs, coefficients that can flip sign. With $R_j^2=0.9$,
    $\mathrm{VIF}_j=1/(1-0.9)=\mathbf{10}$, so $\operatorname{Var}(\hat\beta_j)$ is $10\times$
    larger and its SE is inflated by $\sqrt{10}=\mathbf{3.162}$
    (rule of thumb: $\mathrm{VIF}>5$--$10$ is problematic).
    \smallskip

    \textbf{(e)} Matching (each method once):
    \begin{itemize}
      \item[(i)] \textbf{Linear regression} --- quantitative response; \emph{inference}
        (effect size with CI) from OLS standard errors.
      \item[(ii)] \textbf{Lasso} --- $p=8{,}000\gg n=80$: OLS infeasible; $\ell_1$ selects a
        short list of metabolites.
      \item[(iii)] \textbf{Logistic regression} --- binary response; coefficients give odds
        ratios $e^{\hat\beta_j}$: interpretable for management.
      \item[(iv)] \textbf{Random forest} --- automatic non-linearities and interactions; strong
        accuracy with many predictors; interpretability not needed.
    \end{itemize}
  \end{solutionbox}
\end{frame}

\end{document}
